\documentclass[10pt]{article}
\usepackage[margin=0.9in]{geometry}
\usepackage{amsmath,amssymb,booktabs,enumitem,listings,microtype,xcolor,array}
\usepackage[hidelinks]{hyperref}

\lstdefinelanguage{Patch}{morekeywords={create,number,text,boolean,list,thing,change,called,set,add,remove,clear,show,why,watch,history,undo,redo,preview,if,else,repeat,make,do,return,allow,may,increase,decrease,up,to,true,false,and,or,not,from,window,text,button,input,when,clicked,changed,closed,as},sensitive=true,morecomment=[l]{\#},morestring=[b]"}
\lstset{language=Patch,basicstyle=\ttfamily\small,columns=fullflexible,frame=single,breaklines=true,showstringspaces=false}

\title{Patch: State-Change Factorization and Semantic Change Contracts for Transparent Mutable Programs}
\author{Mich\'el Nguyen}
\date{Beta 8 research artifact manuscript, August 2026}

\begin{document}
\maketitle

\begin{abstract}
General-purpose languages typically expose persistent mutation through assignment or data-structure-specific update operations, while auditing, undo, provenance, replay, and policy enforcement reconstruct additional semantic structure around those writes. Patch explores the opposite factorization: post-creation persistent mutation is represented and executed through a structured semantic change. A programmer writes \texttt{change score: add 1}; the language does not provide ordinary persistent reassignment as an alternative update route. We call the central design property \emph{State-Change Factorization}: every modeled persistent transition from $S$ to $S'$ factors through a semantic change $\delta$ satisfying $apply(\delta,S)=S'$.

Patch derives \emph{Semantic Change Signatures} from this mandatory mutation representation and checks optional \emph{Change Capabilities} constraining target, operation kind, and numeric magnitude. A component may therefore be authorized to increase \texttt{player.score} by at most 10 while remaining unauthorized to replace or decrease the same location. A Lean 4 model proves State-Change Factorization, Mutation Transparency, Change Signature Soundness for a structured control-flow core, and runtime policy containment for protected formal executions.

Beta 8 moves the checked implementation boundary toward source semantics. Generated certificates contain a proof-free formal \texttt{SourceStmt} that preserves source mutation verbs \texttt{add}, \texttt{remove}, \texttt{set}, and \texttt{clear}; a separately produced semantic \texttt{EvidenceStmt}; a separate production Change Signature claim; and the declared policy. Lean validates raw amount bounds and performs source-to-semantic normalization itself. It then checks SourceStmt-to-EvidenceStmt equality, decodes evidence to the formal \texttt{CoreStmt}, reconstructs the formal Change Signature, compares it with the separate production claim, and applies a verified semantic policy checker. For the formal source execution relation, Lean proves that accepted protected source statements cannot emit modeled semantic effects outside policy. The remaining trust gaps are explicit: the JavaScript parser/AST-to-SourceStmt extraction, production interval-analysis soundness, and correspondence between the production runtime and formal execution remain unproved. Patch does not claim effects, capabilities, source calculi, translation validation, proof-carrying evidence, range analysis, provenance, or undo as individually new. The research hypothesis is that mandatory semantic mutation supplies a useful common substrate for quantitative state-transition authority, formal assurance, and runtime tooling without enlarging the beginner-facing language.
\end{abstract}

\section{Introduction}
Mutable programs routinely need answers to questions beyond ``what value is stored now?'' Debuggers need provenance, interfaces need undo, event-sourced systems need replay, and security mechanisms need to constrain effects. In conventional languages these services are often layered around writes whose semantic meaning is only partially represented by the language.

Patch asks what changes if the semantic update is primary. Instead of first mutating state and later attaching metadata, a Patch program describes a change that is itself the route to the new persistent state:

\begin{lstlisting}
create number score = 0

change score:
  add 1
\end{lstlisting}

The surface syntax is deliberately small. Internally the transition can carry target, operation, before/after state, version information, source location, inverse information, provenance, and static authority information.

The same mutation vocabulary enables authority that is more precise than binary write permission:

\begin{lstlisting}
allow reward:
  player.score may increase up to 10

make reward(player, bonus number 0..10):
  change player:
    add bonus to score
\end{lstlisting}

The policy permits a bounded increase. It does not permit an arbitrary \texttt{set}, even though both operations technically write the same path.

The central research question is not whether explicit changes, effects, capabilities, proof-carrying evidence, or verified lowering exist. They do. The question is whether forcing persistent mutation through one semantic change representation gives a useful factorization in which execution, quantitative authority, history, provenance, and formal checking share the same mutation vocabulary.

This beta makes seven artifact-level contributions:
\begin{enumerate}[leftmargin=*]
\item a small language design with no ordinary post-creation persistent-assignment escape hatch;
\item a normalized Change IR and operation-/magnitude-sensitive Change Signatures and Capabilities;
\item a Lean state model proving State-Change Factorization and Mutation Transparency;
\item a structured formal execution model proving Change Signature Soundness and runtime policy containment;
\item a small executable Lean policy checker with machine-checked soundness;
\item a proof-free semantic evidence layer whose decoding and signature correspondence are checked by Lean; and
\item a formal source-core layer preserving \texttt{add}/\texttt{remove}/\texttt{set}/\texttt{clear}, with Lean-side semantic normalization, source/evidence equality checking, source/signature correspondence, and source-level formal policy containment.
\end{enumerate}

The final contribution is intentionally narrower than full compiler verification. JavaScript still parses production Patch source and extracts the formal source representation. Amount ranges are still produced by the JavaScript range analyzer. Beta 8 reduces and exposes these trust boundaries rather than claiming they are solved.

\section{Language Design}
\subsection{Progressive disclosure}
A beginner can ignore signatures, policies, proofs, and certificates:

\begin{lstlisting}
create number lives = 3
change lives:
  remove 1
show lives
\end{lstlisting}

Advanced machinery derives information from the same change model without requiring a second command/event architecture in ordinary user code.

\subsection{No persistent-assignment escape hatch}
\texttt{create} introduces a fresh persistent binding. Existing persistent bindings cannot subsequently be updated using ordinary \texttt{x = e}; state-changing source code must use \texttt{change}. This syntactic restriction aligns the language surface with the single modeled semantic commit route.

\subsection{Source and semantic operation vocabularies}
The source-facing mutation verbs are \texttt{set}, \texttt{add}, \texttt{remove}, and \texttt{clear}. Static semantic analysis can classify numeric mutations more precisely. For example, \texttt{add 5} is an \emph{increase}; \texttt{add -5} is a \emph{decrease}; and \texttt{set 5} is a \emph{set}. Policies can distinguish these operations even when they address the same persistent field.

\subsection{Console and GUI code}
The same source mutation model is used by console and GUI code. For example:

\begin{lstlisting}
create number count = 0

window "Counter":
  text "Count: {count}"
  button "Add" as add_button

when add_button clicked:
  change count:
    add 1
\end{lstlisting}

GUI event execution is not yet part of the formal source subset; it is an implementation case for later correspondence work.

\section{State-Change Factorization}
Let persistent state be a map $\sigma : Name \rightarrow Value$, let $\nu(x)$ be the current version of target $x$, and let history be $H=[\delta_1,\ldots,\delta_n]$.

A normalized semantic change is
\[
\delta = \langle x,b,n,v,\Delta,v'\rangle,
\]
where $x$ is the target, $b$ and $n$ are base and new versions, $v$ and $v'$ are before and after values, and $\Delta$ is an ordered sequence of semantic operations. Well-formedness requires
\[
applyOps(\Delta,v)=v' \quad\text{and}\quad n=b+1.
\]

Define the modeled persistent commit path:
\[
commit(\delta,\sigma,\nu,H)=
\langle\sigma[x\mapsto v'],\nu[x\mapsto n],H\cdot\delta\rangle.
\]

\paragraph{State-Change Factorization.}
For every state-changing step in the current Lean machine, there exists a well-formed $\delta$ such that the resulting machine is exactly the defined commit of that $\delta$.

\paragraph{Mutation Transparency.}
The witnessing change is present in resulting history. Thus a modeled persistent transition cannot occur without an inspectable semantic witness.

These are theorems of the formal machine. Production runtime correspondence remains a separate obligation.

\section{Semantic Change Contracts}
\subsection{Change Signatures}
For a statement or recipe $f$, a Change Signature is a conservative collection of semantic effects:
\[
Sig(f)=\{e_1,\ldots,e_k\}.
\]
An effect identifies a target, optional field, semantic operation kind, and optional amount interval.

For example:

\begin{lstlisting}
make reward(player):
  change player:
    add 5 to score
\end{lstlisting}

is summarized conceptually as:
\begin{verbatim}
player.score -> increase by [5,5]
\end{verbatim}

\subsection{Magnitude-aware capabilities}
A policy may authorize selected semantic transitions:

\begin{lstlisting}
allow reward:
  player.score may increase up to 10
\end{lstlisting}

A protected recipe is rejected by the production compiler if its inferred committed changes are not covered by policy. A \texttt{set} is not interchangeable with an \texttt{increase}, and a provable range whose maximum exceeds 10 is rejected.

Magnitude information alone is not claimed as new. The candidate distinction is that quantitative authority is attached to semantic transition kinds derived from the mandatory mutation representation.

\subsection{Production range reasoning}
Ranged recipe parameters allow the production analyzer to infer bounded dynamic changes:

\begin{lstlisting}
make reward(player, bonus number 0..5):
  change player:
    add bonus * 2 to score
\end{lstlisting}

The current analyzer infers $bonus*2\in[0,10]$. Runtime guards enforce declared parameter bounds. However, the semantic soundness of the JavaScript interval analyzer is not yet mechanized; beta 8 only validates the resulting raw bounds once supplied to the formal source layer.

\section{Lean Core Semantics}
The formal project is pinned to Lean 4.30.0 and currently contains five principal modules.

\paragraph{\texttt{PatchFormal.lean}.}
Defines values, semantic operations, changes, persistent machine state, the modeled change step, intervals, semantic effects, policies, and relational coverage judgments.

\paragraph{\texttt{PatchSignature.lean}.}
Defines the structured semantic core:
\begin{verbatim}
skip | emit effect | seq s1 s2 | branch s1 s2 | repeat n body
\end{verbatim}
with static \texttt{inferSignature} and runtime \texttt{Executes stmt trace}.

\paragraph{\texttt{PatchChecker.lean}.}
Defines executable checks for interval containment, rule/effect admission, policy search, and protected statements, and proves these decisions sound.

\paragraph{\texttt{PatchEvidence.lean}.}
Defines proof-free semantic evidence, validates raw intervals, decodes evidence to \texttt{CoreStmt}, canonicalizes formal signatures, and proves evidence/signature correspondence.

\paragraph{\texttt{PatchSource.lean}.}
Defines the beta-8 formal source mutation vocabulary, source-to-semantic normalization, source/evidence checking, source/signature checking, and a formal source execution relation.

\subsection{Change Signature Soundness}
For the structured semantic core, Lean proves:
\[
Executes(s,r) \Longrightarrow RuntimeChanges(r) \subseteq Sig(s).
\]
Branch inference includes effects from both alternatives, while an execution chooses one. Repetition may emit an effect repeatedly, but every emitted effect remains represented by the static body signature.

\subsection{Capability containment}
If every inferred signature effect is admitted by policy, then:
\[
RuntimeChanges(s) \subseteq Sig(s) \subseteq Cap(s)
\]
and therefore:
\[
RuntimeChanges(s) \subseteq Cap(s).
\]
The formal theorem \texttt{endToEndCapabilitySoundness} establishes this result for the modeled semantic core.

\section{Verified Semantic Policy Checker}
The checker computes policy admission rather than taking it as an unproved premise. Lean proves:
\[
withinBool(I,P)=true \Longrightarrow Within(I,P),
\]
\[
allowsBool(rule,effect)=true \Longrightarrow Allows(rule,effect),
\]
and:
\[
policyAllowsBool(Sig,Policy)=true
\Longrightarrow PolicyAllows(Sig,Policy).
\]

For:
\[
checkProtected(s,P)=policyAllowsBool(inferSignature(s),P),
\]
checker soundness and Change Signature Soundness compose to yield:
\[
\begin{split}
&Executes(s,r) \wedge checkProtected(s,P)=true\\
&\qquad\Longrightarrow \forall e\in r.\;\exists p\in P.\;Allows(p,e).
\end{split}
\]
This theorem is implemented as \texttt{checkedExecutionCannotEscape}.

\section{Proof-Free Semantic Evidence}
Beta 7 removed a producer-generated \texttt{CoreStmt} from the immediate trusted boundary. The producer instead emits:
\begin{verbatim}
EvidenceAmount { lo, hi }
EvidenceEffect { target, field?, kind, amount? }
EvidenceStmt = skip | emit | seq | branch | repeat
\end{verbatim}

An \texttt{EvidenceAmount} contains no proof that $lo\le hi$. \texttt{decodeEvidenceAmount} validates this relation before constructing the semantic \texttt{Interval}. \texttt{decodeEvidenceStmt} recursively constructs the formal \texttt{CoreStmt} only from accepted evidence.

The producer also emits a separate production Change Signature claim. Lean runs the formal \texttt{inferSignature}, canonicalizes it, and compares it to that claim through \texttt{checkEvidenceSignature}. Lean proves:
\[
\begin{split}
&decodeEvidenceStmt(E)=some(s)\\
&\wedge\;checkEvidenceSignature(E,C)=true\\
&\Longrightarrow encodeSignature(inferSignature(s))=C.
\end{split}
\]
The corresponding theorem is \texttt{checkedEvidenceSignatureCorresponds}.

Evidence-level policy checking is also machine checked. \texttt{checkedEvidenceExecutionCannotEscape} composes accepted evidence, formal execution, and the verified policy checker to prove that no modeled runtime effect escapes policy.

\section{Beta 8: Formal Source-Core Correspondence}
\subsection{Why another layer?}
Semantic evidence already improves on directly trusting a producer-generated formal statement, but the producer could still classify a source \texttt{add} as an \texttt{increase} or \texttt{decrease}. Because this semantic distinction is central to Patch policies, beta 8 moves that classification into Lean for the current source subset.

\subsection{Formal source vocabulary}
\texttt{PatchSource.lean} defines:
\begin{verbatim}
SourceChangeKind = add | remove | set | clear
SourceChange { target, field?, kind, rawAmount? }
SourceStmt = skip
           | change SourceChange
           | seq SourceStmt SourceStmt
           | branch SourceStmt SourceStmt
           | repeat Nat SourceStmt
\end{verbatim}

The source representation is proof-free. Numeric \texttt{add}/\texttt{remove} changes carry raw amount bounds supplied by the producer.

\subsection{Lean-side semantic normalization}
The executable \texttt{normalizeSourceChange} validates the raw interval and maps source verbs to semantic effects. For the current fragment:
\begin{center}
\begin{tabular}{ll}
\toprule
Source change & Semantic effect \\
\midrule
\texttt{add [0,5]} & \texttt{increase [0,5]} \\
\texttt{remove [0,5]} & \texttt{decrease [0,5]} \\
\texttt{add [-5,-5]} & \texttt{decrease [5,5]} \\
\texttt{remove [-5,-5]} & \texttt{increase [5,5]} \\
\bottomrule
\end{tabular}
\end{center}

For a non-positive range, mirrored magnitude bounds are passed through the same evidence interval decoder again. Mixed-sign ranges are rejected from the certifiable source fragment because they have no single static semantic direction.

This lets a certificate preserve the distinction between what source said and what semantic evidence claims. For example, source:
\begin{lstlisting}
change player:
  add -5 to score
\end{lstlisting}
remains an \texttt{add} with raw amount $[-5,-5]$ in \texttt{SourceStmt}; the separate semantic evidence claims \texttt{decrease [5,5]}. Lean must derive that relationship.

\subsection{Source-to-evidence equality}
The executable checker:
\begin{verbatim}
checkSourceEvidence source evidence
\end{verbatim}
computes whether Lean's \texttt{lowerSourceStmt source} is exactly the separately supplied evidence. The theorem \texttt{checkSourceEvidence\_sound} establishes:
\[
checkSourceEvidence(S,E)=true
\Longrightarrow lowerSourceStmt(S)=some(E).
\]

Thus semantic classification is no longer accepted solely because the JavaScript producer emitted it.

\subsection{Source-to-signature correspondence}
\texttt{checkSourceSignature} composes source lowering with evidence decoding and formal signature inference:
\begin{verbatim}
SourceStmt
 -> Lean source normalization
 -> EvidenceStmt
 -> CoreStmt
 -> inferSignature
 -> canonical formal signature
 -> compare production signature claim
\end{verbatim}

\texttt{checkSourceSignature\_sound} proves that a successful check yields some evidence and formal statement such that source lowering produces the evidence, evidence decoding produces the statement, and the formal signature exactly equals the separate claim.

\subsection{Formal source execution and policy theorem}
Beta 8 defines:
\begin{verbatim}
SourceExecutes source runtime
\end{verbatim}
by requiring a successful source lowering to evidence, successful evidence decoding to a \texttt{CoreStmt}, and a formal \texttt{Executes} trace with the given runtime semantic effects.

This is intentionally a \emph{formal source-core execution relation}; it is not yet identified with the JavaScript runtime.

Lean proves:
\[
\begin{split}
&SourceExecutes(S,R)\\
&\wedge\;checkSourceProtected(S,P)=true\\
&\Longrightarrow \forall e\in R.\;\exists p\in P.\;Allows(p,e).
\end{split}
\]
The theorem is \texttt{checkedSourceExecutionCannotEscape}. The companion theorem \texttt{checkedSourceSignatureAndPolicy} packages successful source/signature and source/policy checks into one result over the decoded formal statement.

\section{Production Architecture and Certificates}
The production implementation now emits separate views from the real AST:
\begin{verbatim}
                         -> production Change Signature
production Patch AST ---+-> semantic formal bridge / EvidenceStmt
                         `-> formal SourceStmt preserving source verbs
\end{verbatim}

\texttt{src/formal-source.js} performs the source-core extraction. \texttt{src/formal-bridge.js} provides the independent semantic bridge. \texttt{src/change-analysis.js} supplies the production signature claim. The generated certificate therefore does not collapse all claims into one producer representation.

For a protected supported recipe:
\begin{verbatim}
patch certify program.patch --out Program.patchcert.lean
\end{verbatim}
produces:
\begin{itemize}[leftmargin=*]
\item SHA-256 of the exact Patch source bytes;
\item Patch IR, source-schema, and evidence-schema versions;
\item a formal \texttt{SourceStmt};
\item a separate semantic \texttt{EvidenceStmt};
\item a separate production Change Signature claim;
\item the declared semantic policy;
\item a machine-decided SourceStmt-to-EvidenceStmt check;
\item a source-to-formal-signature correspondence theorem;
\item a source-level policy checker theorem; and
\item a formal source-runtime containment theorem.
\end{itemize}

This architecture uses ideas related to translation validation \cite{necula2000translation} and producer/evidence/checker separation from Proof-Carrying Code \cite{necula1997pcc}. Patch claims no novelty for those assurance patterns.

\subsection{Current trusted computing boundary}
Beta 8 still trusts:
\begin{verbatim}
Patch source bytes
  -> JavaScript parser / AST
  -> SourceStmt extraction
\end{verbatim}
and separately trusts that the JavaScript interval analyzer produces a conservative range for each source numeric expression.

After the formal \texttt{SourceStmt} plus its raw amount ranges are supplied, the current Lean chain checks source semantic normalization, semantic evidence equality, formal statement decoding, formal signature reconstruction, production-signature correspondence, semantic policy admission, and runtime containment for the formal source execution relation.

A source SHA-256 binds a certificate to source bytes for artifact identity. It does not prove that the JavaScript frontend extracted the correct SourceStmt.

\section{Implementation and Artifact Status}
The beta-8 production path is:
\begin{verbatim}
Patch source
  -> parser / AST + range annotations
  -> production Change Signature analysis
  -> production capability validation
  -> formal SourceStmt extraction
  -> independent semantic bridge
  -> Change IR 0.6
  -> source/evidence/signature Lean certificate
  -> .patchapp                    [implemented]
  -> bootstrap .wasm              [implemented]
  -> direct Wasm lowering         [next]
  -> native host packaging        [planned]
\end{verbatim}

JavaScript CI runs on Windows, macOS, and Linux with Node 22 and 24. It checks syntax, language behavior, semantic contracts, formal source extraction, semantic bridge behavior, certificate generation, examples, portable bundles, bootstrap Wasm, and the generated project website.

Formal CI explicitly builds \texttt{PatchFormal}, \texttt{PatchSignature}, \texttt{PatchChecker}, \texttt{PatchEvidence}, and \texttt{PatchSource}, then compiles a production-generated certificate. An unfinished-proof gate rejects actual \texttt{sorry}/\texttt{admit} placeholders.

Patch Studio remains a browser-first PWA and project website. It exposes editing, GUI preview, an initial visual Designer, Change Contract inspection, Change IR, and normal program execution. Lean certificate checking currently requires the Node/Lean toolchain.

The current WebAssembly backend emits a valid instantiable module containing Patch source and Change IR for a Patch host. It remains \emph{bootstrap Wasm}; direct Change-IR-to-Wasm execution is not implemented yet.

\section{Related Work and Novelty Boundary}
Patch occupies a crowded research area, so broad priority claims would be inappropriate.

\paragraph{First-class state and state transitions.}
Plaid makes abstract object-state changes first-class \cite{sunshine2011plaid}. Worlds reifies program state and supports scoped speculative side effects and commit/discard \cite{warth2011worlds}. Patch therefore does not claim to invent first-class state transition or reified state.

\paragraph{Effect and capability systems.}
Classical effect systems infer conservative effect approximations \cite{lucassen1988effects}. Effects as Capabilities relates effects and capability-based authority \cite{brachthaeuser2020effects}. Graded, quantitative, refinement, behavioral, permission, and typestate systems provide richer forms of effect/state reasoning. Patch does not claim effect inference, quantitative effects, or effects-plus-capabilities as new. The candidate distinction is specifically that operation- and magnitude-aware state-transition authority is derived from the same semantic change representation through which persistent mutation is required to execute.

\paragraph{Translation validation, source correspondence, and proof-carrying evidence.}
Translation validation checks properties of individual compiler results \cite{necula2000translation}. Proof-Carrying Code separates producer evidence from a smaller verifier \cite{necula1997pcc}. Formal source/intermediate calculi and semantics-preserving lowering are also established compiler-verification techniques. Patch uses these ideas as assurance infrastructure, not novelty claims.

\paragraph{Change calculi and patch systems.}
Edit Lenses model explicit edits across related structures \cite{hofmann2012edit}; higher-order change structures support incremental computation \cite{cai2014changes}; formal patch theories study composition and commutation \cite{angiuli2016patch}. Patch likewise treats changes as semantically structured, but its hypothesis concerns mandatory application-state mutation plus semantic authority.

\paragraph{Provenance and why debugging.}
Whyline demonstrates why/why-not debugging questions \cite{ko2008whyline}. Patch's current \texttt{why} feature is narrower and uses recorded semantic changes. It is supporting tooling, not a primary novelty claim.

\paragraph{Educational representations.}
Scratch demonstrates the value of low-complexity representations and immediate feedback \cite{resnick2009scratch}. Patch's simplicity hypothesis remains empirical; the current formal results do not establish beginner usability.

\section{Evaluation Plan}
A high-venue submission requires evidence beyond artifact unit tests.

\paragraph{E1: Production AST to SourceStmt correspondence.}
Define a stable production-AST subset and independently verify or validate its extraction to the formal \texttt{SourceStmt}. The critical failure criterion is omission, mutation, or invention of a committed source change while still marking the artifact certified.

\paragraph{E2: Interval-analysis soundness.}
Formalize the numeric expression fragment and prove that every concrete evaluation satisfying the range environment lies inside the production-inferred interval. Because magnitude-aware authority is central to the strongest technical claim, this is now the highest-value standalone formal proof.

\paragraph{E3: Production runtime correspondence.}
Relate traces from the actual supported production runtime to the formal \texttt{SourceExecutes}/\texttt{Executes} traces.

\paragraph{E4: Semantic-security case studies.}
Build at least two or three applications in which operation- or magnitude-sensitive authority prevents concrete violations. Candidates include a bounded reward plugin, bounded account debit, and a game/UI extension. At least one case should lie entirely inside the beta-8 source-certified subset.

\paragraph{E5: Conventional baselines.}
Compare Patch against implementations requiring explicit validation, commands/events, logging, or policy wrappers. Measure annotations, security checks, audit plumbing, defect-seeding outcomes, and maintenance burden rather than only source lines.

\paragraph{E6: Assurance cost.}
Measure SourceStmt and EvidenceStmt size, production analysis/extraction time, certificate size, Lean source/evidence equality time, signature-check time, policy-check time, cold/cached verification latency, and scaling with effects and control flow.

\paragraph{E7: Direct compiled execution.}
Implement direct Change-IR-to-WebAssembly execution while preserving a stable correspondence to checked semantic evidence.

\paragraph{E8: Usability, if retained.}
If simplicity remains a headline empirical claim, preregister a controlled study separating beginner mutation comprehension from advanced contract/formal tasks.

\section{Limitations and Threats}
The formal source/core languages remain smaller than production Patch. Recipe calls, dynamic loops, GUI events, undo/redo, richer values, external I/O, and several mutation classes are outside the current source certificate subset. Beta 8 checks correspondence beginning at the emitted formal SourceStmt, not from raw source bytes. The JavaScript parser and SourceStmt extractor are not yet verified.

The production interval analyzer remains unproved against expression evaluation. A raw interval can be internally well ordered and still be semantically unsound if the producer computed it incorrectly. Runtime range guards reduce one operational mismatch but do not replace static soundness. The production runtime is not yet proved to implement the formal SourceExecutes relation.

The formal source execution relation itself is intentionally defined through lowering and decoding into the existing CoreStmt semantics. It therefore establishes a coherent formal pipeline, but should not be described as an independently derived full source operational semantics.

Production call analysis remains conservative and incomplete. Whole before/after snapshots are inefficient for large structures. Provenance-based \texttt{why} is historical rather than counterfactual causation. Bootstrap Wasm is not direct executable lowering.

Novelty remains a threat. Behavioral types, quantitative/refinement effects, capability systems, state-transition languages, and specialized update calculi may contain closer precedents than currently identified. A submission must conduct a systematic literature review and make a precise technical claim rather than rely on the phrase ``change-oriented programming.''

\section{Conclusion}
Patch investigates a specific factorization of mutable programming: persistent mutation is represented and executed as semantic change rather than ordinary assignment with semantic structure reconstructed later. This mandatory mutation substrate supports operation- and magnitude-aware Change Signatures and Capabilities, while Lean proves signature and policy containment for a structured formal core.

Beta 8 moves the assurance boundary toward source semantics. Generated certificates preserve source mutation verbs separately from semantic evidence and production signature claims. Lean validates raw amount bounds, performs source semantic normalization, checks SourceStmt-to-EvidenceStmt equality, reconstructs the formal Change Signature, checks the separate production claim, and applies the verified policy checker. For the formal SourceExecutes relation, accepted protected source statements cannot emit modeled semantic effects outside policy.

The project is not yet a top-venue submission. The strongest next path is deliberately narrow: prove interval-analysis soundness, establish production AST-to-SourceStmt assurance and production-runtime trace correspondence, demonstrate compelling semantic-security cases, measure assurance cost and engineering benefit, and keep the paper centered on mandatory semantic mutation plus quantitative semantic authority rather than the surrounding IDE and product features.

\begin{thebibliography}{11}
\bibitem{sunshine2011plaid} Joshua Sunshine, Karl Naden, Sven Stork, Jonathan Aldrich, and Eric Tanter. First-Class State Change in Plaid. In \emph{OOPSLA}, 2011. \url{https://doi.org/10.1145/2048066.2048122}.

\bibitem{warth2011worlds} Alessandro Warth, Yoshiki Ohshima, Ted Kaehler, and Alan Kay. Worlds: Controlling the Scope of Side Effects. In \emph{ECOOP}, 2011. \url{https://doi.org/10.1007/978-3-642-22655-7_9}.

\bibitem{lucassen1988effects} J. M. Lucassen and D. K. Gifford. Polymorphic Effect Systems. In \emph{POPL}, 1988, pp. 47--57. \url{https://doi.org/10.1145/73560.73564}.

\bibitem{brachthaeuser2020effects} Jonathan Immanuel Brachth\"auser, Philipp Schuster, and Klaus Ostermann. Effects as Capabilities: Effect Handlers and Lightweight Effect Polymorphism. \emph{PACMPL}, 4(OOPSLA), 2020. \url{https://doi.org/10.1145/3428194}.

\bibitem{necula2000translation} George C. Necula. Translation Validation for an Optimizing Compiler. In \emph{PLDI}, 2000, pp. 83--94. \url{https://doi.org/10.1145/349299.349314}.

\bibitem{necula1997pcc} George C. Necula. Proof-Carrying Code. In \emph{POPL}, 1997, pp. 106--119. \url{https://doi.org/10.1145/263699.263712}.

\bibitem{hofmann2012edit} Martin Hofmann, Benjamin C. Pierce, and Daniel Wagner. Edit Lenses. In \emph{POPL}, 2012, pp. 495--508. \url{https://doi.org/10.1145/2103621.2103715}.

\bibitem{cai2014changes} Yufei Cai, Paolo G. Giarrusso, Tillmann Rendel, and Klaus Ostermann. A Theory of Changes for Higher-Order Languages: Incrementalizing Lambda-Calculi by Static Differentiation. In \emph{PLDI}, 2014, pp. 145--155. \url{https://doi.org/10.1145/2594291.2594304}.

\bibitem{angiuli2016patch} Carlo Angiuli, Edward Morehouse, Daniel R. Licata, and Robert Harper. Homotopical Patch Theory. \emph{Journal of Functional Programming}, 26:e18, 2016. \url{https://doi.org/10.1017/S0956796816000198}.

\bibitem{ko2008whyline} Andrew J. Ko and Brad A. Myers. Debugging Reinvented: Asking and Answering Why and Why Not Questions about Program Behavior. In \emph{ICSE}, 2008, pp. 301--310. \url{https://doi.org/10.1145/1368088.1368130}.

\bibitem{resnick2009scratch} Mitchel Resnick et al. Scratch: Programming for All. \emph{Communications of the ACM}, 52(11):60--67, 2009. \url{https://doi.org/10.1145/1592761.1592779}.
\end{thebibliography}

\end{document}
